%قبل از شروع مسئله ترجیحاً چیزی ننویسید. اما اگر خواستید بنویسید، در قالب لیستی (فقرات و شمارش) بنویسید.
%برای اضافه کردن نمره، بعدی از شروع‌مسئله، [x] را اضافه کنید و به جای x  نمره مورد نظر را قرار دهید. برای زیرمسئله‌ها هم به همین صورت است.
%\شروع‌مسئله[10]
\شروع‌مسئله [45] \متن‌سیاه{دیگر داده‌ساختارها}
%\زیرمسئله % مسئلۀ ۱پ
%\زیرمسئله % مسئلۀ ۱ت
\شروع{زیرمسائل}
\زیرمسئله[10] الگوریتمی با استفاده از پشته\پانویس{stack} ارائه دهید که یک آرایه به طول $n$ مثل $A[0,n-1]$ و یک عدد $k$ را از ورودی دریافت کند و کمینۀ هر $k$ عنصر متوالی را در زمان خطی پیدا کند. به طور دقیق‌تر خروجی الگوریتم یک آرایه به طول $n-k+1$ مثل $B[0, n-k]$ است که
$B[i]=\min\{A[i], A[i+1], \dots, A[i+k-1]\}$.
(راهنمایی: عناصر $B$ را یکی یکی حساب کنید. برای استفاده از اطلاعات مراحل قبل در محاسبۀ عنصر $B[i]$ فعلی، باید علاوه بر $B[i-1]$، اطلاعات دیگری را نیز نگه‌داری کنید. چه اطلاعاتی بهتر است نگه‌داری شود و چرا پشته برای نگه‌داری این اطلاعات پشته کفایت می‌کند؟)

\پاسخ{
	روی عناصر آرایه حرکت می‌کنیم و وقتی به عنصر $A[i]$ می‌رسیم، آن را در یک پشته push می‌کنیم با این تفاوت که قبل از push کردن تا جای ممکن اگر عنصر بالای پشته بیشتر از $A[i]$ بود آن را pop می‌کنیم. توجه کنید با این کار از پایین به بالای پشته عناصر مقدارشان صعودی است و اندیس مربوط به آن‌ها نیز صعودی است و بالاترین عنصر در هر لحظه $A[i]$ است. همچنین در هر لحظه اگر $i\geq k-1$ بود، می‌خواهیم کمترین عنصر $A[i-k+1\dots i]$ را خروجی دهیم. برای این کار به پایین‌ترین عنصر پشته نگاه می‌کنیم و اگر خارج از این بازه بود آن را حذف می‌کنیم و در نهایت پایین‌ترین عنصر پشته را به عنوان عنصر کمینه خروجی می‌دهیم. (با توجه به این حذف، می‌توان گفت داده‌ساختار استفاده شده صف دوطرفه است و نه پشته. توجه کنید تأکید بر روی الگوریتم است و توجه به این نکته تأثیری در نمره ندارد.)
	
	در این صورت در هر مرحله یک \متن‌لاتین{push}، تعدادی \متن‌لاتین{pop} و حداکثر یک حذف از پایین پشته داریم. توجه کنید تعداد در مجموع \متن‌لاتین{pop}ها حداکثر $n$ است و بنابراین زمان الگوریتم در مجموع $O(n)$ خواهد بود.
	 
}
%\زیرمسئله آرایه‌ای از اعداد داریم و می‌خواهیم با استفاده از آن یک درخت بسازیم. به این‌صورت که کوچک‌ترین عدد آرایه در ریشه درخت قرار می‌گیرد و زیرآرایه سمت چپ، زیردرخت سمت چپ درخت و زیرآرایه سمت راست، زیردرخت سمت‌راست درخت شود. به وسیله پشته و در زمان خطی این درخت را بسازید.
%این سؤال درخت دکارتی است که در کلاس پوشش داده شد.

\زیرمسئله[5] درخت متناظر عبارت پسوندی زیر را ترسیم کند. سپس عبارت پیشوندی متناظر آن را بنویسید.
$$a\ b + c\ d - *\ e\ f + /$$
%میانوندی زیر را ترسیم کنید و با استفاده از آن نمایش پسوندی و پیشوندی این عبارت را بنویسید. برای این کار از 
%\href{https://fa.wikipedia.org/wiki/اولویت_عملگرها}{جدول اولویت عملگرها در این لینک}
% استفاده کنید.
%$$(a+b) * c \wedge d \wedge e \wedge f * g / h$$

\پاسخ{
$/ * + a\ b - c\ d + e\ f$
	
}

\زیرمسئله[10] تریپ (treap) که ترکیب درخت (tree) و هرم (heap) است، یک درخت دودویی جست‌وجو تصادفی است که ارتفاع آن به احتمال زیاد $\Theta(\log n)$ است. این داده‌ساختار را از کتاب CLRS (ویرایش سوم کتاب، تمرین ۱۳-۴، صفحۀ ۳۳۳) مطالعه کنید و فقط به سؤال قسمت a آن پاسخ دهید.

\پاسخ{
هدف از این سؤال مطالعه مبحث تریپ بوده است و پاسخ قسمت a سؤال نسبتاً ساده است. ریشه $r$ به صورت یکتا تعیین می‌شود زیرا کمترین اولویت را دارد. سپس مجموعه به وسیله کلید ریشه به دو زیرمجموعه یکی با کلیدهای کوچکتر از کلید $r$ و یکی با کلیدهای بزرگ‌تر از کلید $r$ افزار می‌شوند. هر کدام از این زیرمجموعه‌ها نیز یک تریپ تشکیل می‌دهند که زیردرخت‌های چپ و راست $r$ را تشکیل می‌دهند. بنابراین با استفاده از استقرا روی تعداد رأس‌ها می‌توان ثابت کرد که ساختار تریپ به صورت یکتا تعیین می‌شود.
	
}

\زیرمسئله[12] 
از کتاب CLRS (ویرایش سوم) تمرین‌های
17.1-3
(ص ۴۵۶)،
17.2-2
(ص ۴۵۹) و
17.3-2
(ص ۴۶۲)
را حل کنید. سؤال در مورد تحلیل سرشکن یک داده‌ساختار فرضی با سه روش مختلف هست.

\پاسخ{
	روش تجمیعی: مجموع هزینه‌های عملیات برابر است با 
	$$\sum_{i = 1}^n c_i \le \sum_{i = 1}^{n} 1 + \sum_{j = 0}^{\lg n} 2^j < 3n = O(n).$$
	بنابر این هزینۀ سرشکن هر عمل برابر است با $O(n)/n = O(1)$.
	
	روش حسابداری: ادعا می‌کنیم هر عمل هزینۀ سرشکن ۳ دارد. اگر $i$ توانی از ۲ نبود، یک واحد را هزینه می‌کنیم و ۲ واحد را در حساب بانکی ذخیره می‌کنیم. اگر $i$ توانی از ۲ بود، $i$ واحد از حساب بانکی برداشت می‌کنیم. مشابه روش تجمیعی می‌توان نشان داد که حساب بانکی منفی نمی‌شود و بنابراین هزینه سرشکن هر عمل $O(1)$ است.
	
	روش تابع پتانسیل: تابع پتانسل را به صورت
	$\Phi(D_0) = 0$
	و
	 $\Phi(D_i) = 2i - 2^{1 + \lfloor \lg i \rfloor}$ 
برای
	  $i > 0$
 تعریف می‌کنیم.
برای اولین عمل داریم
	$$\hat c_i = c_i + \Phi(D_i) - \Phi(D_{i - 1}) = 1 + 2i - 2^{1+ \lfloor \lg i \rfloor} - 0 = 1.$$
	برای اعمال بعدی اگر $i$ توانی از ۲ نباشد داریم
	$$\hat c_i = c_i + \Phi(D_i) - \Phi(D_{i - 1}) = 1 + 2i - 2^{1 + \lfloor \lg 1 \rfloor} - (2(i - 1) - 2^{1 + \lfloor \lg(i - 1) \rfloor}) = 3.$$
	و اگر $i$ توانی از ۲ باشد داریم
	$$\hat c_i = c_i + \Phi(D_i) - \Phi(D_{i - 1}) = i + 2i - 2^{1 + j}-(2(i - 1) - 2^{1 + j - 1}) = i + 2i - 2i - 2i + 2 + i = 2.$$
	بنابراین هزینه سرشکن هر عمل حداکثر ۳ است.
		
}

\زیرمسئله[8]
با استفاده از آرایۀ پسوندی و آرایۀ کمکی طولانی‌ترین پیشوند مشترک (LCP) برای یک رشتۀ $s$، چطور می‌توان کوچکترین زیررشته ناتهی $s$ را پیدا کرد که فقط یک بار در $s$ تکرار شده است؟

\پاسخ{
تمام خانه‌های آرایۀ پسوندی را به صورت زیر بررسی می‌کنیم. اگر $SA$ آرایه پسوندی باشد و $SA[i]=j$ باشد، در این صورت کوچکترین رشتۀ یکتای $s$ که از $s[j]$ شروع می‌شود را می‌توان با کمک LCP و مقایسه با $s[SA[i-1]]$ و $s[SA[i+1]]$ در $O(1)$ پیدا کرد. در نهایت کمترین مقدار پیدا شده را خروجی می‌دهیم.

}

\پایان{زیرمسائل}





%%%%%%%%%%%%%%%%%%%%%%%%%%%%%%%%%%%%%%%%%%%%%%%%%%%%%
% برخی دستورات مفید برای تمرین‌ها. %
%%%%%%%%%%%%%%%%%%%%%%%%%%%%%%%%%%%%%%%%%%%%%%%%%%%%%

% فرمول‌های نمونه:
%$f(x) = \Theta(x)$
%$T(x, y) \leq \log(x) + 2^y + \binom{2n}{n}$
% {\tt code\_function}


% لیست موارد بدون شماره‌گذاری:
%\شروع{فقرات}
%    \فقره مورد اول 
%        \شروع{فقرات}
%            \فقره مورد اول داخلی 
%                \شروع{فقرات}
%                    \فقره اولین مورد خیلی داخلی 
%                    \فقره دومین مورد خیلی داخلی 
%                \پایان{فقرات}
%            \فقره مورد دوم داخلی 
%        \پایان{فقرات}
%    \فقره مورد دوم 
%\پایان{فقرات}

% لیست موارد شماره‌گذاری شده:
%\شروع{شمارش}
%    \فقره یک مورد 
%    \فقره یک مورد دیگر 
%    \فقره باز هم یک مورد دیگر
%\پایان{شمارش}

% برای نوشتن یک فرمول چند خطی که یک شماره بخوره. برای حالت‌های دیگر به سایت زیر مراجعه کنید:
% https://www.overleaf.com/learn/latex/Aligning_equations_with_amsmath
%\begin{align} 
%    \begin{split}
%        (x+y)^3 &= (x+y)^2(x+y) \\
%                &= (x^2+2xy+y^2)(x+y) \\
%                &= (x^3+2x^2y+xy^2) + (x^2y+2xy^2+y^3) \\
%                &= x^3+3x^2y+3xy^2+y^3
%    \end{split}                                 
%\end{align}

% نوشتن ماتریس و آرایه دوبعدی
%\begin{align}
%    A = 
%    \begin{bmatrix}
%        A_{11} & A_{21} \\
%        A_{21} & A_{22}
%    \end{bmatrix}
%\end{align}

% اضافه کردن تصویر با شرح. برای اضافه کردن تصویر بدون شرح، فقط دستور درج‌تصویر را استفاده کنید.
%\شروع{شکل}[h!]
%\تنظیم‌ازوسط
%  \درج‌تصویر[پهنا=0.5\پهنای‌متن]{img.jpg}
%\شرح{شرح تصویر}
%\برچسب{تصویر:۱}
%\پایان{شکل}
%\پایان{زیرمسائل}

%کد زیر را ببینید.
%%\lstinputlisting{source_filename.py}
%\begin{latin}
%\begin{lstlisting}
%# Import the modules
%import sys
%import random
%
%ans = True
%
%while ans:
%question = raw_input("Ask the magic 8 ball a question: (press enter to quit) ")
%
%answers = random.randint(1,8)
%
%if question == "":
%sys.exit()
%
%elif answers == 1:
%print "It is certain"
%
%elif answers == 2:
%print "Outlook good"
%
%elif answers == 3:
%print "You may rely on it"
%
%elif answers == 4:
%print "Ask again later"
%
%elif answers == 5:
%print "Concentrate and ask again"
%
%elif answers == 6:
%print "Reply hazy, try again"
%
%elif answers == 7:
%print "My reply is no"
%
%elif answers == 8:
%print "My sources say no"
%\end{lstlisting}
%\end{latin}
